% paper_fci_molecules.tex — GeoVac LCAO FCI Paper (Molecules)
% Topological Full Configuration Interaction for Heteronuclear Diatomics
% Status: Revised, March 12, 2026 — corrected for Hamiltonian balance
% Target: arXiv (physics.comp-ph / quant-ph)
\documentclass[aps,pra,twocolumn,superscriptaddress,longbibliography]{revtex4-2}

\usepackage{amsmath,amssymb,amsthm,graphicx,xcolor,bm,braket,dcolumn,booktabs}
\usepackage{hyperref}

\newtheorem{theorem}{Theorem}
\newtheorem{corollary}{Corollary}

\begin{document}

\title{Full Configuration Interaction for Heteronuclear\\
Diatomics in Concatenated Atom-Centred Bases: LiH as a Benchmark}
\thanks{GeoVac FCI Paper in the Geometric Vacuum series}

\author{J.~Loutey}
\affiliation{Independent Researcher, Kent, Washington}

\date{March 12, 2026}

%% ========================================================================
%% ABSTRACT
%% ========================================================================

\begin{abstract}
We present a Topological LCAO architecture for heteronuclear full
configuration interaction (FCI) within the GeoVac discrete graph
framework.  Two independent atom-centered hydrogenic lattices are
concatenated into a bipartite molecular graph via Mulliken cross-nuclear
attraction on the diagonal, Slater--Condon two-electron repulsion
integrals with exact $R^k$ radial integrals and Gaunt coefficients, and
STO-overlap bridges on the off-diagonal; the resulting one- and
two-electron Hamiltonian is diagonalized in the full Slater determinant
space (4\,845 determinants for LiH at $n_{\max}=2$, where the results
below are obtained;\ the solver also handles the 367\,290 determinants of
$n_{\max}=3$).  For LiH the
balanced Hamiltonian produces a monotonically attractive potential
energy surface with no equilibrium geometry; at $n_{\max}=2$ the curve
still lies 0.22~Ha below the separated-atom energy at $R=8$~bohr, so it
does not dissociate to the atoms and no binding energy is extracted.
The absence of a minimum is traced to the $R$-independence of the intra-atom
one-electron block:\ unlike continuous quantum mechanics, each atom's
hydrogenic diagonal $-Z^2/(2n^2)$ is fixed at construction and does not
respond to inter-nuclear compression;\ the large-$R$ offset is not
diagnosed here.  A Fock-weighted kinetic correction, whose
\emph{form} is motivated by the energy-shell constraint $p_0^2 = Z^2/n^2$
of Paper~7 but whose overall scale $\lambda$ is a single fitted parameter
(retained as a labeled qualitative diagnostic, not a zero-parameter
result), partially restores the equilibrium by concentrating
orthogonalization repulsion on compact core orbitals, giving
$R_{\mathrm{eq}}\approx 3.7$~bohr at $\lambda=0.1$ ($n_{\max}=2$)
without repairing the large-$R$ limit.  A 29-version diagnostic arc
(v0.9.8--v0.9.37) systematically tested and eliminated every competing
hypothesis: basis set superposition error as a geometric driver, integral
incompleteness, orbital exponent relaxation, the single-$n$, shared-$p_0$
SO(4) Sturmian bond sphere approximation, and Hamiltonian imbalance between cross-nuclear and
electron repulsion orbital subspaces.  In particular, Wigner $D$-matrix
selection rules $D^{2}_{10,10}(\gamma)=1$ and
$D^{2}_{00,10}(\gamma)=0$ are proven as structural results.  The
excitation-driven Direct CI solver achieves
$O(N_{\mathrm{SD}}\times N_{\mathrm{connected}})$ scaling, preserving the
$O(V)$ character of the underlying graph Hamiltonian.
\end{abstract}

\maketitle

%% ========================================================================
\section{Introduction}
\label{sec:intro}
%% ========================================================================

The GeoVac program works on the discrete label set $(n,l,m)$ of the
unit three-sphere $S^3$, which carries the single-particle
quantum-number structure without any continuous spatial
integration~\cite{paper0,paper1,paper7}.  Prior work established the
exact correspondence at the level of labels for single-center atomic
systems---on them the $S^3$ Laplacian is the diagonal $SO(4)$ Casimir,
whose spectrum Fock's projection maps to the hydrogenic one---with $O(V)$
scaling;\ the graph Laplacian itself does not converge to the $S^3$
operator (Paper~7, corrected 2026-09-25).  Extending this framework to multi-center molecules, however,
presents a fundamental geometric challenge: two or more atomic lattices,
each living on its own $S^3$ conformal manifold, must be bridged into a
single molecular graph without resorting to the continuous
$\mathbb{R}^3$ integration that the discrete premise is designed to
replace.

In this work we introduce the \emph{Topological LCAO} architecture,
which solves the multi-center problem by concatenating atomic subgraphs
into a bipartite molecular graph connected by discrete off-diagonal
bridges.  The resulting block Hamiltonian is fed into an $O(V)$
excitation-driven Full Configuration Interaction solver that operates
entirely in Slater determinant space.  The headline result is negative:
for LiH the balanced Hamiltonian gives a monotonically attractive
potential energy surface with no equilibrium geometry, revealing a
fundamental limitation of the discrete LCAO framework:\ the absence of
$R$-dependent kinetic repulsion.  A Fock-weighted kinetic correction,
whose form follows Paper~7~\cite{paper7} but whose scale $\lambda$ is
fitted, partially restores an equilibrium.

LiH is the minimal four-electron heteronuclear diatomic and therefore
the ideal first benchmark for the Topological LCAO framework.  Its four
electrons permit exact FCI treatment without truncation of the many-body
expansion, while the $Z_A=3$, $Z_B=1$ asymmetry breaks the inversion
symmetry exploited by homonuclear codes and exposes every heteronuclear
coupling pathway.  The experimental targets are well-established:
$D_e = 0.092$~Ha and $R_{\mathrm{eq}}=3.015$~bohr~\cite{huber}.  This
paper formally extends the companion computational framework described
in Paper~8~\cite{paper8}.

The paper is organized as follows.  Section~\ref{sec:architecture}
develops the bipartite graph construction and the $O(V)$ Direct CI
solver.  Section~\ref{sec:results} presents the LiH benchmark results,
including the counterpoise-corrected potential energy surface and the
Hamiltonian balance requirement.
Section~\ref{sec:diagnostic} documents the 29-version diagnostic arc
that systematically falsified every hypothesis for the observed
behavior---integral incompleteness, orbital
relaxation, and Sturmian geometries.
Section~\ref{sec:kinetic} identifies the root cause as missing
$R$-dependent kinetic repulsion and presents a Fock-weighted correction
whose form follows Paper~7 and whose scale is fitted.
Section~\ref{sec:conclusion} summarizes and outlines next steps.

%% ========================================================================
\section{Topological LCAO Architecture}
\label{sec:architecture}
%% ========================================================================

%% ------------------------------------------------------------------------
\subsection{Subgraph Concatenation}
\label{sec:subgraph}
%% ------------------------------------------------------------------------

The \texttt{MolecularLatticeIndex} instantiates two independent atomic
subgraphs: a lithium lattice with nuclear charge $Z_A=3$ and a hydrogen
lattice with $Z_B=1$.  Each subgraph is a scale-invariant hydrogenic
lattice defined on its own $S^3$ conformal manifold~\cite{paper7},
retaining its atom-centered conformal structure; the hydrogenic energies
$\varepsilon_{n} = -Z^2/(2n^2)$ are assigned to the labels $n$ through
Fock's energy shell $p_0^2 = -2E$ (the bare graph Laplacian is
positive-semidefinite and does not converge to $\Delta_{S^3}$, so these
negative energies are neither its eigenvalues nor a continuum image of
its spectrum).  At $n_{\max}=3$, the lithium
subgraph contributes $\sum_{n=1}^{3}n^2 = 14$ spatial orbitals and the
hydrogen subgraph contributes another 14, yielding a combined molecular
basis of 28 spatial orbitals (56 spin-orbitals).

The concatenation is a direct sum at the graph level: the adjacency
matrix of the molecular graph is block-diagonal in the two atomic
subgraphs prior to bridging.  No spatial overlap integral is computed
between the subgraphs at this stage; the topological identity of each
atomic lattice is preserved exactly.

%% ------------------------------------------------------------------------
\subsection{Topological Bridging and the Hamiltonian}
\label{sec:bridging}
%% ------------------------------------------------------------------------

Bonding emerges from a rank-update concatenation that introduces
off-diagonal coupling between the two atomic blocks.  The one-electron
molecular Hamiltonian takes the block form
\begin{equation}
\mathbf{H}_1 =
\begin{pmatrix}
\mathbf{H}_A & \mathbf{V}_{AB} \\
\mathbf{V}_{AB}^{\dagger} & \mathbf{H}_B
\end{pmatrix},
\label{eq:block_h1}
\end{equation}
where $\mathbf{H}_A$ and $\mathbf{H}_B$ are the isolated atomic
Hamiltonians and $\mathbf{V}_{AB}$ encodes the inter-atomic coupling.

The diagonal blocks carry the hydrogenic eigenvalues
$\varepsilon_n = -Z^2/(2n^2)$ (the energy-shell map of the labels $n$,
not graph-Laplacian eigenvalues and not an image of the graph spectrum) supplemented
by the Mulliken
cross-nuclear attraction: each orbital $\phi_i^{(A)}$ centered on atom
$A$ experiences the potential of nucleus $B$ via the diagonal matrix
element
\begin{equation}
\langle \phi_i^{(A)} | {-Z_B}/{r_B} | \phi_i^{(A)} \rangle,
\label{eq:mulliken}
\end{equation}
evaluated analytically from the hydrogenic Fourier representation.
Off-diagonal elements of $\mathbf{V}_{AB}$ are STO-overlap bridges
weighted by conformal factors derived from the stereographic
projection~\cite{fock1935}.

The two-electron repulsion $V_{ee}$ is constructed via Slater--Condon
rules with two classes of integrals.  Same-atom Coulomb and exchange
integrals use exact Slater $R^k$ radial integrals with Gaunt angular
coefficients, validated against the known analytical values
$F^0(1s,1s)=5Z/8$, $F^0(1s,2s)=17Z/81$, and
$F^0(2s,2s)=77Z/512$.  Cross-atom Coulomb repulsion
is evaluated via multipole expansion, with exchange restricted to
$s$-orbital pairs where the overlap is non-negligible.

%% ------------------------------------------------------------------------
\subsection{\texorpdfstring{$O(V)$}{O(V)} Direct Configuration Interaction}
\label{sec:direct_ci}
%% ------------------------------------------------------------------------

Distributing four electrons across 56 spin-orbitals generates
$\binom{56}{4} = 367\,290$ Slater determinants.  A na\"ive pairwise
construction of the CI Hamiltonian matrix would require
$O(N_{\mathrm{SD}}^2)$ evaluations---over $10^{11}$ matrix element
computations---rendering exact diagonalization impractical.

We instead adopt the excitation-driven Direct CI algorithm of Knowles
and Handy~\cite{knowles_handy}, which constructs the Hamiltonian action
$\sigma = \mathbf{H}|\Psi\rangle$ by traversing only the nonzero edges
generated by valid single and double excitations from each determinant.
The computational cost scales as
$O(N_{\mathrm{SD}} \times N_{\mathrm{connected}})$, where
$N_{\mathrm{connected}}$ is the average number of determinants connected
to each reference by at most a double excitation.  For the LiH basis at
$n_{\max}=3$, this reduces wall time from hours to approximately
8~seconds when the Slater--Condon assembly loops are JIT-compiled
via Numba~\cite{numba} (approximately 75~seconds in pure Python).

The one-electron integrals $h_{pq}$ and two-electron integrals
$(pq|rs)$ are stored as dense arrays (56 spin-orbitals yield lookup
tables of order $56^2$ and $28^4$, both fitting comfortably in memory),
eliminating the $\sim\!24\,\mu$s per-access overhead of
\texttt{scipy.sparse} index validation that dominated earlier
implementations.

%% ------------------------------------------------------------------------
\subsection{Counterpoise Correction Framework}
\label{sec:cp}
%% ------------------------------------------------------------------------

Basis set superposition error (BSSE) is an inherent artifact of any
atom-centered expansion: at finite basis size, each atom can
variationally borrow orbitals from its partner, artificially lowering
the molecular energy below the true complete-basis-set limit.  We
implement the Boys--Bernardi counterpoise (CP)
correction~\cite{boys_bernardi} natively within the Topological LCAO
framework.

Ghost fragments are constructed by preserving the full bipartite
topology of the molecular graph while zeroing the nuclear charge and
electron population of the partner center.  Specifically, a ghost-$B$
calculation retains all 28 spatial orbitals but sets $Z_B=0$,
eliminating all bridges, cross-nuclear attraction, and cross-atom
$V_{ee}$ involving center $B$.  The fragment energy
$E_A^{\mathrm{mol}}$ is then computed in the identical topological basis
space as the supermolecule, guaranteeing a consistent subtraction:
\begin{equation}
D_e^{\mathrm{CP}} = E_A^{\mathrm{mol}} + E_B^{\mathrm{mol}} -
E_{AB}(R_{\mathrm{eq}}).
\label{eq:cp}
\end{equation}

%% ========================================================================
\section{LiH Benchmark Results}
\label{sec:results}
%% ========================================================================

%% ------------------------------------------------------------------------
\subsection{Computational Parameters}
\label{sec:params}
%% ------------------------------------------------------------------------

The molecular graph is assembled from two atomic subgraphs at
$n_{\max}=3$: a lithium lattice ($Z_A=3$, 14 spatial orbitals) and a
hydrogen lattice ($Z_B=1$, 14 spatial orbitals), yielding a combined
basis of 28 spatial orbitals and 56 spin-orbitals.  The full
four-electron Slater determinant space contains $\binom{56}{4} =
367\,290$ configurations.  The exact ground-state eigenvalue is obtained
from the excitation-driven Direct CI solver described in
Sec.~\ref{sec:direct_ci}, with wall time of approximately 8~seconds on
a single CPU core (Numba-accelerated assembly; approximately 75~seconds
without JIT compilation).  The energetics below are reported at
$n_{\max}=2$ (5 spatial orbitals per atom, 20 spin-orbitals, 4\,845
determinants), the size at which the physical one-electron Hamiltonian
has been run; the $n_{\max}=3$ figures above describe the method and
its cost.

%% ------------------------------------------------------------------------
\subsection{Energetic Accuracy and Basis Set Superposition}
\label{sec:energy}
%% ------------------------------------------------------------------------

Table~\ref{tab:lih_energies} gives the LiH potential energy surface at
$n_{\max}=2$ with the physical one-electron Hamiltonian: the hydrogenic
diagonal plus cross-nuclear attraction, with the STO-overlap bridges as
the only off-diagonal one-electron terms.  The separated-atom reference
(each atom in its own basis) is $E_{\mathrm{Li}}+E_{\mathrm{H}} =
-7.601$~Ha, and the Boys--Bernardi counterpoise correction gives
$\Delta E_{\mathrm{BSSE}} = 0.088$~Ha, the same at every $R$ from 1.5 to
8~bohr (the ghost fragments carry no $R$-dependent term).  The surface does not return to
the separated-atom energy at large $R$ (Sec.~\ref{sec:contraction}), so
the energy difference to the atomic fragments at $R=3.015$~bohr is not a
dissociation energy and is not compared with experiment.

\begin{table}[h]
\centering
\begin{tabular}{cc}
\toprule
$R$ (bohr) & $E$ (Ha) \\
\midrule
1.5 & $-8.3793$ \\
2.5 & $-8.0592$ \\
3.015 & $-7.9899$ \\
4.0 & $-7.9312$ \\
5.0 & $-7.8877$ \\
8.0 & $-7.8258$ \\
\midrule
separated atoms & $-7.601$ \\
\bottomrule
\end{tabular}
\caption{LiH potential energy surface at $n_{\max}=2$ (balanced
\texttt{exact+True} configuration, physical one-electron Hamiltonian).
The energy falls monotonically as $R$ decreases, with no minimum, and at
$R=8$~bohr it lies 0.22~Ha below the separated-atom energy.}
\label{tab:lih_energies}
\end{table}

\paragraph{Hamiltonian balance requirement.}
The cross-nuclear attraction and electron repulsion must operate on
matched orbital subspaces.  Two balanced configurations---\texttt{exact}
cross-nuclear with full cross-atom $V_{ee}$ (\texttt{exact+True}), and
\texttt{fourier} cross-nuclear with $s$-only $V_{ee}$
(\texttt{fourier+s\_only})---produce nearly identical molecular energies
across the PES.  An \emph{imbalanced} configuration
(\texttt{exact+s\_only}) artificially overbinds, as the $p$ and $d$
orbitals receive cross-nuclear attraction without the compensating
electron repulsion screening.  The results presented here are from the
balanced \texttt{exact+True} Hamiltonian.

%% ------------------------------------------------------------------------
\subsection{Absence of Equilibrium Geometry}
\label{sec:contraction}
%% ------------------------------------------------------------------------

The balanced configurations (\texttt{exact+True} and
\texttt{fourier+s\_only}) both yield monotonically attractive PES curves
with no equilibrium geometry.

With the physical one-electron Hamiltonian (Table~\ref{tab:lih_energies})
the energy falls from $-7.826$~Ha at $R=8.0$~bohr to $-8.379$~Ha at
$R=1.5$~bohr, the slope diminishing at large $R$ but never reversing
sign, and at $R=8$~bohr the curve still lies 0.22~Ha below the
separated-atom energy.  At no point does a repulsive wall balance the
attractive cross-nuclear and exchange interactions: the system is
variationally attracted to short bond lengths at all $R$.

This monotonic behavior is traced to a structural property of the
Topological LCAO: the one-electron block of each atom is
\emph{completely $R$-independent}.  In continuous quantum
mechanics, bringing two atoms together forces their electrons into
higher-momentum orthogonal states, creating the kinetic repulsion that
establishes the equilibrium bond length.  In the discrete framework,
each atom's one-electron block---the hydrogenic diagonal $-Z^2/(2n^2)$,
which carries that atom's kinetic energy---is fixed at construction time
and does not respond to the approach of a second center:\ the
atom-centered orbitals are treated as orthonormal across the two
centers, so no orthogonalization kinetic energy arises.

The only $R$-dependent off-diagonal one-electron terms are the bridge
edges connecting the two subgraphs, and they are small:\ the maximum
bridge matrix element at $R=3.015$~bohr is 0.017~Ha, and the bridges
produce no repulsive wall.

This absence of equilibrium geometry motivates both the systematic
diagnostic arc in the following section and the kinetic repulsion
analysis in Sec.~\ref{sec:kinetic}.

%% ========================================================================
\section{Systematic Diagnostic Arc}
\label{sec:diagnostic}
%% ========================================================================

The initial implementation (v0.9.8) exhibited heavy overbinding
($D_e^{\mathrm{raw}} = 0.205$~Ha, more than twice the experimental
value) and a 17\% geometric contraction of the equilibrium bond length.
Over the course of 29 controlled experiments spanning versions v0.9.8
through v0.9.36, every plausible hypothesis for these discrepancies was
systematically tested.  The experiments are grouped below by hypothesis
class; the complete version history is tabulated in
Appendix~\ref{app:versions}.

%% ------------------------------------------------------------------------
\subsection{Hypothesis 1: Basis Set Superposition Error
(v0.9.8--v0.9.11)}
\label{sec:hyp_bsse}
%% ------------------------------------------------------------------------

The first hypothesis attributed both the energetic overestimation and
the geometric contraction to BSSE: electrons borrowing orbitals from the
partner center would artificially lower the molecular energy, inflating
$D_e$ and potentially distorting the PES minimum.

The Boys--Bernardi counterpoise correction (Sec.~\ref{sec:cp})
quantified the superposition error
($\Delta E_{\mathrm{BSSE}} = 0.088$~Ha at $n_{\max}=2$ with the
physical one-electron Hamiltonian, Table~\ref{tab:lih_energies}).
BSSE does not address the absence of an equilibrium geometry, which
has a distinct root cause (Sec.~\ref{sec:kinetic}).

%% ------------------------------------------------------------------------
\subsection{Hypothesis 2: Integral Incompleteness
(v0.9.12--v0.9.13, v0.9.35)}
\label{sec:hyp_integrals}
%% ------------------------------------------------------------------------

The second hypothesis class attributed the geometric contraction to
missing or approximate molecular integrals: either the Mulliken
cross-nuclear attraction overestimates the short-range potential, or the
restriction of cross-atom $V_{ee}$ to $s$-orbital pairs neglects
significant contributions from higher angular momentum channels.

Exact cross-nuclear attraction matrix elements, computed analytically
from the hydrogenic Fourier representation and validated against the
shell theorem, deviate by less than $2$~mHa from the Mulliken
diagonal approximation at $n_{\max}=2$ ($R=2.5$--$5.0$~bohr),
confirming that the Mulliken scheme is close to exact for this system.

With the Fourier cross-nuclear diagonal, extending the cross-atom
$V_{ee}$ to all $(n,l)$ pairs---392 Coulomb integrals at $n_{\max}=3$
compared to 18 for the $s$-only restriction---shifts the total energy
by only $+0.3$ to $+0.4$~mHa at $n_{\max}=2$.  This negligible shift reflects the fact that the LiH
ground state is overwhelmingly $s$-orbital dominated: the $1s^2 2s^2$
configuration of lithium and the $1s$ of hydrogen leave higher angular
momentum channels essentially unoccupied.

Integral incompleteness is conclusively falsified as a mechanism for the
geometric contraction.

%% ------------------------------------------------------------------------
\subsection{Hypothesis 3: Orbital Exponent Relaxation (v0.9.36)}
\label{sec:hyp_relaxation}
%% ------------------------------------------------------------------------

The third hypothesis proposed that fixing atomic orbital exponents at
their free-atom values forces artificial proximity: if electrons cannot
polarize toward the bonding region by adjusting their radial extent, the
system compensates by contracting $R_{\mathrm{eq}}$.

A variational orbital exponent parameter $\zeta_B$ was introduced for
the hydrogen center, allowing the effective nuclear charge
$Z_{\mathrm{eff}} = \zeta_B Z_B$ to optimize the radial extent of all
hydrogen-centered orbitals.  At $n_{\max}=2$ with the physical
one-electron Hamiltonian (scan step $0.1$), the variational minimum
occurs at $\zeta_B \approx 1.3$ at $R=3.015$~bohr, relaxing toward
unity at large $R$ ($\zeta_B \approx 1.1$ at $R=5$ and $8$~bohr).

The relaxation energy is approximately $-0.31$~Ha, but it is
\emph{uniform across the bonding region}: the energy lowering at
$R=2.5$~bohr and $R=3.015$~bohr differs by only $0.004$~Ha.
Consequently the relaxed PES remains monotonically attractive, with no
minimum.  Moreover, the energy below the counterpoise ghost fragments at
$R=3.015$~bohr grows by a factor of $2.2$ (from $0.300$ to
$0.652$~Ha), demonstrating that the variational freedom is exploited
for unphysical basis compression rather than physical polarization.

Orbital exponent relaxation is eliminated as a mechanism for geometric
contraction.

%% ------------------------------------------------------------------------
\subsection{Hypothesis 4: Bond Sphere and Sturmian Basis
(v0.9.15--v0.9.34)}
\label{sec:hyp_sturmian}
%% ------------------------------------------------------------------------

The fourth and most extensive hypothesis class proposed that the
flat-space LCAO concatenation is fundamentally incompatible with the
$S^3$ topology, and that a proper molecular treatment requires Sturmian
basis functions defined on a single \emph{bond sphere}---an $S^3$
manifold with two weighted poles---connected by SO(4) Wigner $D$-matrix
elements.

This program was pursued across 20 versions (v0.9.15--v0.9.34),
progressing from the initial SO(4) Wigner $D$-matrix implementation (44
tests, v0.9.15) through geometric-mean coupling (v0.9.16, unbound),
Shibuya--Wulfman off-diagonal nuclear attraction (v0.9.17, unbound),
hybrid Fourier--SW schemes (v0.9.18, 55\% overbinding), atom-centered
Sturmian variants (v0.9.20--v0.9.28, all unbound or massively
overbinding), prolate spheroidal MO projections
(v0.9.29--v0.9.30, unphysical), full MO Sturmian FCI
(v0.9.31--v0.9.33, 12--25\% error), and dual-$p_0$ extensions (v0.9.34,
25\% error with $p_0$ values collapsing to a single geometric mean).

The definitive negative result is the structural theorem proven in
companion Paper~8~\cite{paper8} (Theorem~4): in a Sturmian basis
sharing a single momentum parameter $p_0$, treated with Paper~8's
single-$n$ rotation approximation, the one-electron Hamiltonian
and overlap obey the two-term identity
$H_{ij} = -\tfrac12 p_0^2 S_{ij} + (1-\beta_j)V_{ij}$ and share a single
orthogonal $SO(4)$ Wigner-$D$ congruence $D^{(n)}(\gamma(R))$ (so $H$ is
block-, not scalar-, proportional to $S$).  The generalized eigenvalues
are therefore independent of the bond angle $\gamma$, hence of $R$,
yielding pure $1/R$ nuclear repulsion with no electronic binding.  Within that
approximation, recovering $R$-dependent binding energies requires solving the continuous
two-center prolate spheroidal PDE for the spectrum $\beta_k(R)$,
reintroducing exactly the continuous spatial geometry that the discrete
graph framework is designed to eliminate.  The Bond Sphere Sturmian
alternative is therefore falsified \emph{in this single-$n$, shared-$p_0$
form}.  The scope matters:\ the genuine cross-$n$ Coulomb--Sturmian overlap
couples $n$ and escapes the theorem, so binding is not excluded for it
(Paper~8, Remark).  What is
ruled out is this approximation, not the discrete framework.

\paragraph{Cross-$n$ hybridization (v0.9.37).}
A further hypothesis attributed the geometric error to the $D$-matrix's
block-diagonality in $n$: the single-$n$ $D$-matrix approximation of
Paper~8 connects only same-$n$ orbitals (Li~$1s \leftrightarrow$ H~$1s$,
Li~$2s \leftrightarrow$ H~$2s$), giving zero off-diagonal coupling for
the chemically relevant Li~$2s \leftrightarrow$ H~$1s$ pair, which the
full cross-$n$ Shibuya--Wulfman matrix does couple~\cite{paper8}.  Overlap-based bridges using the
Shibuya--Wulfman energy scale $-(Z_{\mathrm{eff}}/p_0)\sin\gamma$ were
implemented, producing Li~$2s$--H~$1s$ coupling of 0.052~Ha (40\% of
the same-$n$ $D$-matrix coupling).  However, the PES worsened:\ the
cross-$n$ coupling adds binding at short~$R$, where the overlap is
largest.  Cross-shell hybridization is eliminated as a
mechanism for the geometry error.

\paragraph{Level~4 locked-core (v0.9.37).}
The Level~4 molecule-frame hyperspherical solver (Paper~15) was extended
to LiH by locking the Li~$1s^2$ core and treating the two active
electrons in hyperspherical coordinates with effective charges
$Z_{\mathrm{eff}}$.  Four core models were tested (bare $Z$, integer
screening, Clementi--Raimondi, and Core Hartree with penetration); all
produce PES that collapse to short bond lengths
($R_{\mathrm{eq}} \lesssim 2.5$~bohr) with $D_e > 12$~eV.  The failure has the same root cause:
the core-to-H-nucleus attraction $V_{\mathrm{cross}} = -2Z_B\langle
1s^2_Z | 1/r_B | 1s^2_Z \rangle$ behaves as $-2/R$ at large $R$ but
saturates exponentially at short $R$ (exp$(-6R)$ decay for Li $1s$),
providing insufficient repulsion to establish a minimum at
$R \approx 3$~bohr.

\medskip

Having falsified BSSE as a geometric driver, eliminated integral
incompleteness, ruled out orbital exponent relaxation, shown that the
single-$n$, shared-$p_0$ Sturmian approximation cannot bind (Paper~8's
theorem, which does not exclude the genuine cross-$n$ basis), and
excluded cross-$n$ hybridization and locked-core hyperspherical
approaches, we trace the absence of an equilibrium geometry in the
balanced Hamiltonian to the $R$-independence of the intra-atom
one-electron block, analyzed in the following section.

%% ========================================================================
\section{Missing Kinetic Repulsion and the Fock-Weighted Correction}
\label{sec:kinetic}
%% ========================================================================

%% ------------------------------------------------------------------------
\subsection{The Kinetic Diagnostic}
\label{sec:kinetic_diagnostic}
%% ------------------------------------------------------------------------

Direct examination confirms that the intra-atom one-electron block (the
hydrogenic diagonal $-Z^2/(2n^2)$ within each atomic subgraph) is
identically $R$-independent: it is fixed by the isolated-atom labels and
never modified as $R$ changes.  Bridge edges contribute a small
$R$-dependent off-diagonal term, but the maximum bridge element at
$R=3.015$~bohr is only 0.017~Ha, and it produces no repulsive wall.

This is a structural limitation of the Topological LCAO architecture:
the concatenated basis, treated as orthonormal across the two centers,
has no mechanism to represent the orthogonalization kinetic energy that
arises when atomic orbitals overlap in continuous quantum mechanics.

%% ------------------------------------------------------------------------
\subsection{The Fock-Weighted Kinetic Correction}
\label{sec:fock_correction}
%% ------------------------------------------------------------------------

The energy-shell constraint of Paper~7---$p_0^2 = -2E = Z^2/n^2$ for
the $n$-th shell---provides a natural momentum scale $p_0 = Z/n$ for
each orbital.  We exploit this to construct a physically motivated
kinetic correction.

For each orbital $\phi_a$ on atom $A$, the Fock-weighted overlap
kinetic correction adds to the diagonal of $\mathbf{H}_1$:
\begin{equation}
\Delta h_{aa} = \lambda \sum_{b \in B}
S_{ab}^2 \left(\frac{Z_A}{n_a}\right)^2
\left(\frac{Z_B}{n_b}\right)^2,
\label{eq:fock_correction}
\end{equation}
where $S_{ab} = \langle \phi_a | \phi_b \rangle$ is the STO overlap
integral between $s$-orbitals on different centers, and $\lambda$ is a
single adjustable parameter.  The Fock weights suppress diffuse
orbitals (high $n$) and concentrate the kinetic penalty on compact core
orbitals whose overlap decays rapidly with $R$.

The physical content of Eq.~(\ref{eq:fock_correction}) is that the cost
of overlapping two orbitals scales with the product of their squared
shell momenta $p_0^2 = Z^2/n^2$.  Core--core
pairs (Li~1$s$--H~1$s$) carry weight $9 \times 1 = 9$, while
diffuse--diffuse pairs (Li~3$s$--H~3$s$) carry weight
$1 \times 0.11 = 0.11$---an 81$\times$ discrimination.

The $R$-discrimination (ratio of correction at $R=1.5$ to $R=5.0$~bohr)
improves from 3.0$\times$ with uniform weighting to 9.3$\times$ with
Fock weighting.  For the Li~1$s$ orbital alone, the ratio improves from
19$\times$ to 74$\times$.

%% ------------------------------------------------------------------------
\subsection{Results of the Fock-Weighted Correction}
\label{sec:fock_results}
%% ------------------------------------------------------------------------

Table~\ref{tab:fock_scan} summarizes the $\lambda$-scan at
$n_{\max}=2$ with the physical one-electron Hamiltonian
($R = 1.5$--$8.0$~bohr, step $0.5$~bohr near the minimum).

\begin{table}[h]
\centering
\begin{tabular}{lccc}
\toprule
$\lambda$ & Equilibrium? & $R_{\mathrm{eq}}$ (bohr) &
$E(8.0)-E(R_{\mathrm{eq}})$ (Ha) \\
\midrule
0.00 & No & --- & --- \\
0.02 & No & --- & --- \\
0.05 & No & --- & --- \\
\textbf{0.10} & \textbf{Yes} & $\bm{\sim 3.7}$ & \textbf{0.092} \\
0.20 & Yes & $\sim 4.1$ & 0.075 \\
0.50 & Yes & $\sim 4.7$ & 0.045 \\
\bottomrule
\end{tabular}
\caption{Fock-weighted kinetic correction scan at $n_{\max}=2$
($\lambda$ fitted; $R_{\mathrm{eq}}$ from a three-point parabola about
the lowest grid point).  For $\lambda \geq 0.1$ a potential energy
minimum appears and moves outward as $\lambda$ grows; the last column
is the depth below the $R=8$~bohr point, not a dissociation energy.}
\label{tab:fock_scan}
\end{table}

For $\lambda \geq 0.1$ the Fock-weighted correction creates a potential
energy minimum with a repulsive inner wall, at
$R_{\mathrm{eq}} \approx 3.7$~bohr for $\lambda = 0.1$, moving to
$4.1$ and $4.7$~bohr at $\lambda = 0.2$ and $0.5$.  This is a
qualitative success---the first equilibrium geometry produced by any
balanced Hamiltonian configuration---but it does not repair the large-$R$
limit:\ for every $\lambda$ the curve at $R=8$~bohr still lies
$0.22$~Ha below the separated-atom energy, so no dissociation energy is
obtained, and no scanned value of $\lambda$ places the minimum at the
experimental $3.015$~bohr.

%% ------------------------------------------------------------------------
\subsection{Remaining Limitations}
\label{sec:limitations}
%% ------------------------------------------------------------------------

The remaining discrepancy ($R_{\mathrm{eq}} \approx 3.7$ vs.\
experiment $3.015$~bohr at $\lambda=0.1$, and no return to the
separated atoms) likely reflects the
limitation of a purely diagonal correction.  In standard quantum
chemistry, kinetic repulsion includes both diagonal (orbital compression)
and off-diagonal (orthogonalization) contributions from the full kinetic
energy matrix $T_{AB} = \langle\phi_A|-\tfrac{1}{2}\nabla^2|\phi_B\rangle$.
The Fock-weighted correction captures only the diagonal component.

The BSSE of $0.088$~Ha at $n_{\max}=2$ is comparable to the
experimental dissociation energy of $0.092$~Ha, confirming significant
basis deficiency.  However, since the balanced PES has no
minimum, BSSE is not the mechanism for geometric contraction---it is a
separate issue of basis completeness that CP correction addresses for
vertical binding energies.

%% ========================================================================
\section{Conclusion}
\label{sec:conclusion}
%% ========================================================================

We have presented a heteronuclear full configuration interaction
calculation in concatenated atom-centered hydrogenic bases on the $S^3$
label set.  The Topological LCAO
architecture concatenates atom-centered hydrogenic lattices via Mulliken
cross-nuclear attraction, Slater--Condon two-electron integrals, and
STO-overlap bridges, and diagonalizes the resulting Hamiltonian in the
full Slater-determinant space using excitation-driven Direct CI
(4\,845 determinants at $n_{\max}=2$, where the LiH results are
obtained;\ the solver also handles the 367\,290 determinants of
$n_{\max}=3$).

For LiH the balanced Hamiltonian produces a monotonically attractive
PES with no equilibrium geometry; at $n_{\max}=2$ it still lies 0.22~Ha
below the separated-atom energy at $R=8$~bohr.  The absence of a
minimum is traced to the $R$-independence of each atom's one-electron
block, a structural property rather than basis truncation;\ the
large-$R$ offset is not diagnosed here.

A 29-version diagnostic arc (v0.9.8--v0.9.37) systematically falsified
every competing hypothesis: basis set superposition error as a geometric
driver, integral incompleteness, orbital exponent relaxation, and the
single-$n$, shared-$p_0$ SO(4) Sturmian Bond Sphere approximation.  The
missing kinetic repulsion was identified as a structural limitation of
the LCAO concatenation, and a Fock-weighted kinetic correction, whose
form follows Paper~7's energy-shell constraint $p_0^2 = Z^2/n^2$ and
whose scale $\lambda$ is fitted, partially restores the
equilibrium ($R_{\mathrm{eq}} \approx 3.7$~bohr at $\lambda = 0.1$,
$n_{\max}=2$) without repairing the large-$R$ limit.

The $O(V)$ scaling of the underlying graph Hamiltonian is preserved
throughout, with the Direct CI solver achieving
$O(N_{\mathrm{SD}}\times N_{\mathrm{connected}})$ assembly cost.
The natural next steps are:
\begin{enumerate}
\item Extending the Fock-weighted correction to include off-diagonal
kinetic matrix elements between atomic centers, which would provide
the missing intermediate-$R$ binding;
\item Increasing $n_{\max}$ to assess basis convergence of the
potential energy surface and its large-$R$ limit.
\item \textbf{Prolate spheroidal lattice:} The fundamental limitation
identified above---that the LCAO one-electron blocks provide no
$R$-dependent kinetic repulsion---arises because the molecular graph
is constructed as a concatenation of atomic $S^3$ subgraphs whose
one-electron blocks are fixed at construction time.  This limitation
is resolved by recognizing that the natural geometry for two-center
systems is the prolate spheroid, not the concatenated $S^3$.  In
prolate spheroidal coordinates $(\xi, \eta, \varphi)$, the two-center
Coulomb problem separates exactly, and working in the label set and
channel structure of these coordinates yields a molecular lattice
with $R$-dependence built into the metric.  For H$_2^+$, this approach
achieves $R_{\mathrm{eq}} = 2.001$~bohr (0.21\% error) at the
finite-difference level, while the spectral Laguerre radial solver
reaches machine precision, both with zero fitted
parameters~\cite{loutey_paper11}.  The LCAO framework documented in
this paper remains valuable for understanding the anatomy of molecular
binding in a discrete framework, and the diagnostic arc
(Table~\ref{tab:versions}) provides a systematic record of approaches
that fail and why.
\end{enumerate}
The Fock projection's energy-shell scale $p_0^2 = Z^2/n^2$ supplies a
qualitatively useful weighting for the kinetic repulsion, with its
overall scale $\lambda$ fitted.

%% ========================================================================
\appendix
\section{Version History}
\label{app:versions}
%% ========================================================================

Table~\ref{tab:versions} documents the complete diagnostic arc from
v0.9.11 (the balanced LCAO baseline) through v0.9.38, including
architecture, results, and root cause classification for each version.

\begin{table*}[t]
\centering
\small
\begin{tabular}{llccll}
\toprule
Version & Architecture & $D_e^{\mathrm{CP}}$ (Ha) & $R_{\mathrm{eq}}$ (bohr)
& Status & Root Cause \\
\midrule
v0.9.11 & LCAO baseline (balanced \texttt{exact+True})
  & --- & No eq. & Baseline & No large-$R$ return to atoms \\
v0.9.12 & L\"owdin orthogonalization
  & --- & --- & Fail & Breaks variational principle \\
v0.9.13 & Mulliken exchange $K$
  & --- & $\sim$2.5 & --- & $K$ overcorrects $D_e$ \\
v0.9.15 & SO(4) Wigner $D$-matrix
  & --- & --- & Infrastructure & 44/44 tests, no production \\
v0.9.16 & $D$-matrix cross-nuclear (geom.\ mean)
  & --- & --- & Unbound & Coupling too weak \\
v0.9.17 & Shibuya--Wulfman off-diagonal
  & --- & --- & Unbound & SW coupling too weak \\
v0.9.18 & Hybrid (Fourier diag + SW off-diag)
  & 0.143 & $<$2.0 & 55\% err & Dissociation limit wrong \\
v0.9.20 & Sturmian diagonal
  & --- & --- & Unbound & 1$s$ diagonal positive \\
v0.9.21 & $p_0$-consistent $D$-matrix + self-loop
  & --- & --- & Unbound & Compactness overbinding \\
v0.9.22 & Atom-dependent $p_0$
  & 3.777 & $\sim$2.0 & $3.8\times$ overbind & Scales do not decouple \\
v0.9.23 & Shell-radius attenuation
  & --- & --- & Negligible & Correction $<$ 1 mHa \\
v0.9.24 & Energy decomposition diagnostic
  & --- & --- & Tool & Identifies Fourier diagonal \\
v0.9.25 & Ohno--Klopman cross-atom $V_{ee}$
  & --- & --- & Negligible & Cross-atom $J$ irrelevant \\
v0.9.26 & Angular-weighted $V_{ee}$
  & --- & --- & Negligible & $\Delta E \sim 0.4$ mHa \\
v0.9.27 & SO(4) CG cross-center $V_{ee}$
  & --- & --- & Unbound & $J_{\mathrm{cross}} > J_{\mathrm{same}}$ \\
v0.9.28 & Bond sphere node weights
  & --- & --- & Unbound & $\beta_H$ 46\% short \\
v0.9.29 & Prolate spheroidal MO betas
  & --- & --- & Unbound & LiH mapping ad hoc \\
v0.9.30 & MO-to-atom projection
  & --- & --- & Unphysical & $E=-42.5$ Ha ($5.3\times$) \\
v0.9.31 & MO Sturmian FCI (full)
  & --- & --- & Unbound/overbinding & Shared $SO(4)$ congruence, eigenvalues $\gamma$-independent \\
v0.9.32 & MO Sturmian FCI (variant)
  & --- & --- & Unbound/overbinding & Shared $SO(4)$ congruence, eigenvalues $\gamma$-independent \\
v0.9.33 & MO Sturmian FCI (variant)
  & --- & --- & Unbound/overbinding & Shared $SO(4)$ congruence, eigenvalues $\gamma$-independent \\
v0.9.34 & Dual-$p_0$ extension
  & --- & --- & Unbound/overbinding & Shared $SO(4)$ congruence, eigenvalues $\gamma$-independent \\
v0.9.35 & All-$l$ cross-atom $V_{ee}$
  & --- & No eq. & $\Delta E = 0.4$~mHa & Negligible correction \\
v0.9.36 & Orbital exponent relaxation
  & --- & No eq. & Fail & Uniform PES shift \\
v0.9.37 & Hamiltonian balance diagnostic
  & --- & No eq. & Corrected & Imbalance identified \\
v0.9.38 & Numba-accelerated assembly
  & --- & No eq. & $738\times$ speedup & Python loop overhead eliminated \\
\bottomrule
\end{tabular}
\caption{Diagnostic arc version history (v0.9.11--v0.9.38).
The LCAO baseline (v0.9.11, balanced \texttt{exact+True}) lacks an
equilibrium geometry due to missing kinetic repulsion and does not
return to the separated-atom energy at large $R$
(Table~\ref{tab:lih_energies}).}
\label{tab:versions}
\end{table*}

%% ========================================================================
%% REFERENCES
%% ========================================================================

\begin{thebibliography}{9}

\bibitem{paper0}
J.~Loutey, ``Counting Angular Momentum Labels by Isotropic Packing,'' GeoVac Paper~0 (2026).

\bibitem{paper1}
J.~Loutey, ``The Geometric Atom: the Hydrogen State Space as a Label Lattice,''
GeoVac Paper~1 (2026).

\bibitem{paper7}
J.~Loutey, ``The Dimensionless Vacuum: Recovering the Schr\"odinger
Equation on a Scale-Invariant Label Lattice,'' GeoVac Paper~7 (2026).

\bibitem{paper8}
J.~Loutey, ``Bond Sphere: Diatomic Molecules as $S^3$ and their Natural
Sturmian Basis,'' GeoVac Papers~8--9 (2026).

\bibitem{fock1935}
V.~Fock, ``Zur Theorie des Wasserstoffatoms,''
Z.~Phys. \textbf{98}, 145 (1935).

\bibitem{boys_bernardi}
S.~F.~Boys and F.~Bernardi, ``The calculation of small molecular
interactions by the differences of separate total energies.  Some
procedures with reduced errors,''
Mol.~Phys. \textbf{19}, 553 (1970).

\bibitem{huber}
K.~P.~Huber and G.~Herzberg, \emph{Molecular Spectra and Molecular
Structure IV: Constants of Diatomic Molecules}
(Van Nostrand Reinhold, New York, 1979).

\bibitem{knowles_handy}
P.~J.~Knowles and N.~C.~Handy, ``A new determinant-based full
configuration interaction method,''
Chem.~Phys.~Lett. \textbf{111}, 315 (1984).

\bibitem{numba}
S.~K.~Lam, A.~Pitrou, and S.~Seibert, ``Numba: A LLVM-based Python
JIT compiler,'' in \emph{Proceedings of the Second Workshop on the LLVM
Compiler Infrastructure in HPC} (ACM, 2015), pp.~1--6.

\bibitem{loutey_paper11}
J.~Loutey,
``The Molecular Fock Projection: Diatomic Molecules as Prolate
Spheroidal Lattices,''
GeoVac Paper~11 (2026).

\end{thebibliography}

\end{document}
